\subsection{58. PINN-RANS: Physics-Informed Neural Networks for RANS Equations (Eivazi et al.\ 2022)}
\textbf{arXiv:} \texttt{2107.10711}\quad\textbf{Authors:} Eivazi, Tahani, Schlatter, Vinuesa (2022)\quad\textbf{Journal:} Physics of Fluids 34, 075117\quad\textbf{Domain:} ML / PDE / CFD\\
\scorebadge{Coverage}{6}\quad\scorebadge{Agreement}{5}\quad\textbf{Total:} 11/20\par\smallskip
\noindent\textbf{Coverage rationale.} 3 of 5 test cases attempted (Falkner-Skan BL, ZPG TBL, periodic hill); 2 blocked by unavailable DNS/LES reference datasets (APG TBL from Bobke et al.\ 2017, NACA4412 from Vinuesa et al.\ 2018). No public code from the authors. 14/25 claims tested (56\%); 11/25 data-blocked (100\% accounted for). Architecture reimplemented in PyTorch (paper used TensorFlow): 8 hidden layers $\times$ 20 neurons, tanh activation, Adam$\to$L-BFGS training.\par
\noindent\textbf{Agreement rationale.} All 5 qualitative claims confirmed (PINN-RANS framework works; boundary-data + PDE-residual training converges; pressure predicted without pressure BC). However, 0/14 quantitative claims reproduced within tolerance: streamwise velocity errors 4--67$\times$ higher than paper (e.g., FSBL $E_U$: 4.68\% vs paper's 0.07\%; ZPG $E_{uv}$: 73.6\% vs paper's 6.46\%). Reynolds stress errors in periodic hill are catastrophic (100--13500$\times$), primarily due to synthetic reference data replacing unavailable DNS.\par\smallskip
\noindent\textbf{Key findings:}
\begin{itemize}[leftmargin=1.4em,itemsep=0.2em,topsep=0.2em]
\item The PINN-RANS concept is sound: PINNs can solve RANS without a turbulence model
\item Quantitative reproduction fails without access to the original DNS/LES datasets
\item Even the laminar Falkner-Skan case (analytical solution) shows 67$\times$ higher error, suggesting L-BFGS implementation and hyperparameter sensitivity as contributing factors
\item Framework difference (TensorFlow$\to$PyTorch) and L-BFGS implementation differences are secondary contributors
\end{itemize}
\noindent\textbf{Verdict: PARTIAL.} Methodology confirmed; quantitative claims not reproduced.
\noindent\textbf{Follow-on research questions:}
\begin{enumerate}[leftmargin=1.6em,itemsep=0.2em,topsep=0.2em,label=Q\arabic*.]
\item Contact the KTH-FlowAI group for the DNS datasets and re-run---does access to the correct reference data close the quantitative gap?
\item Implement NTK-based adaptive loss weighting for the multi-term PINN-RANS loss---does this reduce the sensitivity to loss balance coefficients?
\item Compare TensorFlow and PyTorch L-BFGS implementations on a controlled PINN benchmark (e.g., 1D Poisson) to isolate framework-specific convergence differences.
\item How does the PINN-RANS approach scale with Reynolds number? Test at Re$_{\theta}$=500, 2000, 5000, 10000 to identify where the method breaks down.
\item Can transfer learning from the Falkner-Skan solution improve PINN convergence for the turbulent cases? Pre-train on laminar, fine-tune on turbulent BCs.
\end{enumerate}

